% !TeX root = ../../book.tex
\subsection*{集合}

\begin{chapex}
用指定的表示法表示以下集合。
\begin{enumerate}[(a)]
\item 用列表表示法表示 $\{ n \in \mathbb{Z} \mid n^2 < 20 \}$；
\item 用隐含列表表示法表示$\{ 4k+3 \mid k \in \mathbb{N} \}$；
\item 用集合构造表示法表示 6 的所有奇数倍数的集合；
\item 用集合构造表示法表示集合 $\{ 1, 2, 5, 10, 17, \dots, n^2+1, \dots \}$。
\end{enumerate}
\end{chapex}

\begin{chapex}
对于每个 $n \in \mathbb{N}$ 找到集合 $X_n$，满足对于所有 $n \in \mathbb{N}$, $X_{n+1} \subsetneqq X_n$。$X_n$ 集合中可以为空吗？
\end{chapex}

\begin{chapex}
用列表表示法表示集合 $\mathcal{P}(\{ \varnothing, \{ \varnothing, \{ \varnothing \} \} \})$。
\end{chapex}

\begin{chapex}
设 $X$ 为集合，设 $U,V\subseteq X$。证明 $U$ 和 $V$ 不相交当且仅当 $U \subseteq X \setminus V$。
\end{chapex}

\begin{chapex}
对于以下每个陈述，判断是否对所有集合 $A$ 和 $X$ 都成立，并证明你的判断。
\begin{multicols}{2}
\begin{enumerate}[(a)]
\item 如果 $X \setminus A = \varnothing$，则 $X = A$。
\item 如果 $X \setminus A = X$，则 $A = \varnothing$。
\item 如果 $X \setminus A = A$，则 $A = \varnothing$。
\item $X \setminus (X \setminus A) = A$。
\end{enumerate}
\end{multicols}
\end{chapex}

\begin{chapex}
\label{cqPowerSetRespectSetOperations}
对于以下每个陈述，判断是否对于所有集合 $X,Y$ 都成立，还是都不成立，或是某些情况成立，某些情况不成立。
\begin{multicols}{2}
\begin{enumerate}[(a)]
\item $\mathcal{P}(X \cup Y) = \mathcal{P}(X) \cup \mathcal{P}(Y)$
\item $\mathcal{P}(X \cap Y) = \mathcal{P}(X) \cap \mathcal{P}(Y)$
\item $\mathcal{P}(X \times Y) = \mathcal{P}(X) \times \mathcal{P}(Y)$
\item $\mathcal{P}(X \setminus Y) = \mathcal{P}(X) \setminus \mathcal{P}(Y)$
\end{enumerate}
\end{multicols}
\end{chapex}

\begin{chapex}
设 $F$ 为集合，其元素都是集合。证明如果 $\forall A \in F,~ \forall x \in A,~ x \in F$，则 $F \subseteq \mathcal{P}(F)$。
\end{chapex}

\Crefrange{cqSymmetricDifferenceBegin}{cqSymmetricDifferenceEnd} 涉及集合的 \textit{对称差}，定义如下。

\begin{definition}
\label{defSymmetricDifference}
\index{symmetric difference}
\nindex{symmdiff}{$\triangle$}{symmetric difference}
集合 $X$ 和 $Y$ 的 \textbf{对称差} 是集合 $X \symmdiff Y$ \inlatex{triangle}\lindexmmc{triangle}{$\triangle$} 定义为
\[ X \symmdiff Y = \{ a \mid a \in X \text{ 或 } a \in Y \text{ 但不同时属于两个集合} \} \]
\end{definition}

\begin{chapex}
\label{cqSymmetricDifferenceBegin}
证明对于所有集合 $X$ 和 $Y$, $X \symmdiff Y = (X \setminus Y) \cup (Y \setminus X) = (X \cup Y) \setminus (X \cap Y)$。
\end{chapex}

\begin{chapex}
设 $X$ 为集合。证明 $X \symmdiff X = \varnothing$ 和 $X \symmdiff \varnothing = X$。
\end{chapex}

\begin{chapex}
设 $X$ 和 $Y$ 为集合。证明 $X = Y$ 当且仅当 $X \symmdiff Y = \varnothing$。
\end{chapex}

\begin{chapex}
证明集合 $X$ 和集合 $Y$ 不相交当且仅当 $X \symmdiff Y = X \cup Y$。
\end{chapex}

\begin{chapex}
证明对于所有集合 $X$, $Y$ 和 $Z$, $X \symmdiff (Y \symmdiff Z) = (X \symmdiff Y) \symmdiff Z$。
\hintlabel{cqSymmetricDifferenceAssociative}{%
写出一长串方程似乎很诱人，但通过双重包含来证明这一点要容易得多，并在需要时分成不同的情况。一个更简单的方法是巧妙地使用真值表。
}
\end{chapex}

\begin{chapex}
\label{cqSymmetricDifferenceEnd}
证明对于所有集合 $X$, $Y$ 和 $Z$, $X \cap (Y \symmdiff Z) = (X \cap Y) \symmdiff (X \cap Z)$。
\hintlabel{cqIntersectionDistributesOverSymmetricDifference}{%
\Cref{cqSymmetricDifferenceAssociative} 的提示也适用于此。
}
\end{chapex}

\begin{definition}
\label{defOpenSubsetOfR}
\index{open!subset of $\mathbb{R}$}
设集合 $U \subseteq \mathbb{R}$，如果对于所有 $a \in U$，存在 $\delta > 0$ 使得 $(a-\delta, a+\delta) \subseteq U$，则称 $U$ 是\textbf{开放的}。
\end{definition}

在问题 \Crefrange{cqOpenSubsetsOfRBegin}{cqOpenSubsetsOfREnd} 中，您需要证明有关 $\mathbb{R}$ 开子集的一些基本事实。

\begin{chapex}
\label{cqOpenSubsetsOfRBegin}
对于 $\mathbb{R}$ 的以下每个子集，判断（并证明）是否是开放的：
\begin{multicols}{3}
\begin{enumerate}[(a)]
\item $\varnothing$;
\item $(0,1)$;
\item $(0,1]$;
\item $\mathbb{Z}$;
\item $\mathbb{R} \setminus \mathbb{Z}$;
\item $\mathbb{Q}$.
\end{enumerate}
\end{multicols}
\end{chapex}

\begin{chapex}
证明子集 $U \subseteq \mathbb{R}$ 为开集，当且仅当对于所有 $a \in U$，存在 $u,v \in \mathbb{R}$ 使得 $u<a <v$ 且 $(u,v) \subseteq U$。
\end{chapex}

\begin{chapex}
在本题中，您将证明有限多个开集的交集是开集，但无限多个开集的交集可能不是开集。
\begin{enumerate}[(a)]
\item 设 $n \ge 1$ 并假设 $U_1, U_2, \dots, U_n$ 是 $\mathbb{R}$ 的开子集。证明交集 $U_1 \cap U_2 \cap \cdots \cap U_n$ 为开集。
\item 证明 $(0,1+\frac{1}{n})$ 对所有 $n \ge 1$ 开放，但 $\bigcap_{n \ge 1} (0,1+\textstyle\frac{ 1}{n})$ 不开放。
\end{enumerate}
\end{chapex}

\begin{chapex}
\label{cqOpenSubsetsOfREnd}
证明子集 $U \subseteq \mathbb{R}$ 为开集当且仅当它可以表示为开区间的并集 --- 更准确地说，$U \subseteq \mathbb{R}$ 为开集当且仅当对于某个索引集 $I$，对于每个 $i \in I$ 都存在实数 $a_i, b_i$，使得 $U = \bigcup_{i \in I} (a_i, b_i)$。
\end{chapex}

\begin{chapex}
设 $\{ A_n \mid n \in \mathbb{N} \}$ 和 $\{ B_n \mid n \in \mathbb{N} \}$ 为集合族，满足对于所有 $i \in \mathbb{N}$，存在某个 $j \ge i$ 使得 $B_j \subseteq A_i$。证明 $\bigcap_{n \in \mathbb{N}} A_n = \bigcap_{n \in \mathbb{N}} B_n$。
\end{chapex}

\subsection*{判断题}

\tfquestiontext{cqSetsTFBegin}{cqSetsTFEnd}

\begin{chapex}
\label{cqSetsTFBegin}
如果 $E$ 为集合且 $\neg \forall x,\, x \in E$，则 $E = \varnothing$。
\end{chapex}

\begin{chapex}
如果 $E$ 为集合且 $\forall x,\, \neg x \in E$，则 $E = \varnothing$。
\end{chapex}

\begin{chapex}
$\{ 1, 2, 3 \} = \{ 1, 2, 1, 3, 2, 1 \}$
\end{chapex}

\begin{chapex}
$\varnothing \in \varnothing$
\end{chapex}

\begin{chapex}
$\varnothing \in \{ \varnothing \}$
\end{chapex}

\begin{chapex}
\label{cqSetsTFEnd}
$\varnothing \in \{ \{ \varnothing \} \}$.
\end{chapex}

\subsection*{可能性问题}

\asnquestiontext{cqSetsASNBegin}{cqSetsASNEnd}

\begin{chapex} % Always
\label{cqSetsASNBegin}
设 $a$ 和 $b$ 为任意对象。则 $\{ a, b \} = \{ b, a \}$。
\end{chapex}

\begin{chapex} % Sometimes
设 $a$, $b$ 和 $c$ 则 $\{ a, b \} = \{ a, c \}$。
\end{chapex}

\begin{chapex} % Sometimes
设 $X$ 和 $Y$ 为集合，假设存在某个 $a$ 满足 $a \in X \Leftrightarrow a \in Y$。则 $X = Y$。
\end{chapex}

\begin{chapex} % Never
设 $X$ 和 $Y$ 为集合，假设存在某个 $a$ 满足 $a \in X$ 但 $a \not\in Y$，则 $X = Y$。
\end{chapex}

\begin{chapex} % Sometimes
设 $X$ 和 $Y$ 为集合，且 $X \ne Y$。则 $\forall a,\, a \in X \Rightarrow a \not\in Y$。
\end{chapex}

\begin{chapex} % Always
设 $X$ 和 $Y$ 为集合，且 $X \cap Y = \varnothing$。则 $\forall a,\, a \in X \Rightarrow a \not\in Y$。
\end{chapex}

\begin{chapex} % Always
设 $X$ 和 $Y$ 为集合，且 $X \ne Y$。则 $\exists a,\, a \in X \wedge a \not\in Y$。
\end{chapex}

\begin{chapex} % Sometimes
设 $X$ 和 $Y$ 为集合的集合，并假设 $X \in Y$。则 $X \subseteq Y$。
\end{chapex}

\begin{chapex} % Always
设 $X$ 为集合。则 $\varnothing \in \mathcal{P}(X)$。
\end{chapex}

\begin{chapex} % Sometimes
设 $X$ 为集合。则 $\{ \varnothing \} \in \mathcal{P}(X)$。
\end{chapex}

\begin{chapex} % Always
设 $X$ 为集合。则 $\{ \varnothing, X \} \subseteq \mathcal{P}(X)$。
\end{chapex}

\begin{chapex} % Always
设 $X$ 和 $Y$ 为集合，且满足 $\mathcal{P}(X) = \mathcal{P}(Y)$。则 $X = Y$。
\end{chapex}

\begin{chapex} % Sometimes
设 $X$ 和 $Y$ 为集合，且满足 $X \setminus Y = \varnothing$。则 $X = Y$。
\end{chapex}

\begin{chapex} % Always
设 $X$, $Y$ 和 $Z$ 为集合，满足 $X \setminus Y = Z$ 且 $Y \subseteq Z$。则 $X = Y \cup Z$。
\end{chapex}

\begin{chapex} % Sometimes
设 $X$ 和 $Y$ 为集合。则 $X \cup Y = X \cap Y$。
\end{chapex}

\begin{chapex} % Sometimes
\label{cqSetsASNEnd}
设 $X$ 为集合，并设 $A \subseteq X$。则 $A \in X$。
\end{chapex}